% Self-contained source using the official ICLR 2027 style.
% The current method is called NBPO throughout; earlier fixed-reference results are labeled as a precursor variant.
\documentclass{article}
\usepackage{iclr2027_conference,times}
\usepackage[hidelinks,hypertexnames=false]{hyperref}
\usepackage{url}
\usepackage{graphicx}
\usepackage{amsmath}
\usepackage{amssymb}
\usepackage{amsthm}
\usepackage{booktabs}
\usepackage{algorithm}
\usepackage{algorithmic}
\usepackage{float}
\usepackage{microtype}
\usepackage{placeins}

\newtheorem{theorem}{Theorem}
\newtheorem{lemma}{Lemma}
\newtheorem{proposition}{Proposition}
\newtheorem{corollary}{Corollary}
\newtheorem{assumption}{Assumption}
\newtheorem{definition}{Definition}
\theoremstyle{remark}
\newtheorem{remark}{Remark}
\theoremstyle{plain}

\newcommand{\Y}{\mathcal{Y}}
\newcommand{\G}{\mathcal{G}}
\newcommand{\KL}{\mathrm{KL}}
\newcommand{\E}{\mathbb{E}}
\newcommand{\R}{\mathbb{R}}
\newcommand{\D}{\mathcal{D}}
\newcommand{\piref}{\pi_{\mathrm{ref}}}
\newcommand{\NBPO}{\textsc{NBPO}}
\newcommand{\pending}{\textemdash}
\newcommand{\meanstd}[2]{\ensuremath{#1{\scriptstyle\pm#2}}}
% The experiment runner may overwrite this path before compilation.
\providecommand{\NBPOResultDir}{experiments/iclr2027_table1_v2/results}

\title{Nash Bargaining Preference Optimization}
\author{Anonymous authors}
%\iclrfinalcopy

\begin{document}
\maketitle

\begin{abstract}
Aligning one large language model (LLM) to several objectives requires two decisions: how to represent preferences within each objective and how to choose one compromise across objectives. The first decision is difficult because pairwise judgments can be cyclic in realistic settings. For a medical-advice prompt, annotators may prefer a concise answer to a verbose one, the verbose answer to an incomplete one, yet the incomplete answer to the concise one because it includes a critical safety caveat. No scalar ranking can satisfy all three comparisons. Nash Learning from Human Feedback avoids this restriction by treating pairwise preferences as game payoffs, but its multi-objective extensions still require an aggregation rule. We introduce Nash Bargaining Preference Optimization (\NBPO{}), which evaluates each objective through a regularized preference game and applies Nash bargaining as an aggregation rule to its improvement over a common reference policy. Our main technical contribution is a practical optimizer for these policy-dependent game values: a constrained KL-proximal update handles zero surplus at the reference, closed-form opponent reweighting exposes weak responses, and pairwise log-ratio regression removes the intractable policy normalizer while learning directly from binary comparisons. Under exact population updates, the final iterate achieves convergence with rate $O(1/T)$. Experiments compare preference-based and reward-model-based multi-objective baselines under independent evaluation, while controlled rescaling tests verify the predicted invariance.
\end{abstract}

\section{Introduction}
\label{sec:intro}
Reinforcement Learning from Human Feedback (RLHF) is a standard approach for aligning large language models (LLMs) with human judgments \citep{christiano2017deep,ouyang2022training}. Standard RLHF learns a scalar reward from pairwise comparisons. Direct Preference Optimization removes the separately trained reward model, but still assumes that comparisons are generated by an implicit scalar reward through a Bradley--Terry (BT) model \citep{rafailov2023direct,bradley1952rank}. Both approaches therefore require every response to admit one scalar score whose ordering explains all pairwise preferences.

This assumption can fail even within a single alignment objective. Human judgments often reflect several latent attributes whose relative importance changes with the prompt and the pair being compared. Consider three answers to a medical-advice question. Response $A$ is concise and actionable but omits an important caveat; response $B$ is comprehensive but verbose; and response $C$ is less complete but includes the critical safety warning. Annotators may prefer $A$ to $B$ for relevance, $B$ to $C$ for completeness, and $C$ to $A$ for safety and reliability. These comparisons are cyclic: $A\succ B$, $B\succ C$, and $C\succ A$. Fitting the first two comparisons with a scalar reward requires $r(A)>r(B)>r(C)$, which contradicts the third. As a result, optimizing the fitted scalar reward can favor a response that appears best under the learned ranking but is still beaten by another response according to the original pairwise judgments.

Nash Learning from Human Feedback (NLHF) avoids this misspecification by treating prompt-conditioned comparisons directly as payoffs in a zero-sum preference game, and INPO, SPPO, and related methods provide scalable equilibrium-seeking updates \citep{munos2024nash,zhang2025iterative,wu2025sppo,calandriello2024online,rosset2024direct,zhou2025egpo}. Extending this view to multiple objectives introduces a second problem: after representing each objective without imposing transitivity, which compromise should one deployed policy realize? Scalar multi-objective methods expose this trade-off through weights, priorities, constraints, policy interpolation, or preference conditioning \citep{zhou2024beyond,rame2023rewarded,zhong2024panacea,chakraborty2024maxmin,son2026rmod,agnihotri2026mopo}. Recent game-theoretic methods preserve general pairwise feedback but combine objectives through Blackwell target sets or a worst-objective criterion \citep{blackwell1956analog,bhatia2020multicriteria,zhang2026prosper}.

Classical Nash bargaining offers a priority-free way to select one compromise. It maximizes the product of each objective's improvement over a common fallback \citep{nash1950}. In alignment, the reference policy provides a natural fallback: an aligned policy should improve each objective relative to the model that would otherwise be deployed. The classical Nash solution is Pareto optimal and invariant to independent changes in the units used for different objectives. Prior alignment work has applied Nash welfare to learned scalar rewards, while game-based preference methods have focused mainly on minimax or Blackwell solutions \citep{zhong2024nashwelfare,navon2022nashmtl,zhang2026prosper}. What remains is a practical method for Nash bargaining when every objective value is itself defined by a policy-dependent preference game rather than by a fixed scalar reward.

We introduce Nash Bargaining Preference Optimization (\NBPO{}) for this setting. Within each objective, \NBPO{} evaluates the learner against an adaptive opponent that emphasizes responses exposing the learner's weaknesses, while a KL penalty keeps the opponent near the reference policy. Across objectives, it maximizes the Nash product of the resulting improvements over that same reference. The standard bargaining properties are useful in LLM alignment because they ensure that no objective falls below the fallback at the population optimum, no other feasible policy can improve one objective without worsening another, and arbitrary judge-specific units do not alter the selected compromise.

The main technical challenge is optimization, not re-deriving the classical bargaining axioms. Three issues arise simultaneously: the reference policy has zero improvement and lies on the boundary of the logarithmic objective; each adaptive opponent depends on the policy being optimized; and the ideal exponential policy update contains an intractable normalizing constant. We address these issues with a constrained KL-proximal formulation, closed-form opponent reweighting, a low-dimensional dual update whose weights are inversely proportional to achieved improvements, and partition-free pairwise log-ratio regression from binary comparisons.

Our contributions are:
\begin{itemize}
    \item We formulate multi-objective alignment as Nash bargaining over objective-wise regularized preference-game values, thereby separating within-objective preference representation from across-objective compromise selection.
    \item We derive a practical training procedure that handles the zero-surplus reference boundary, alternates closed-form opponent updates with policy fitting, and learns the policy update without fitting scalar rewards or evaluating a language-model partition function.
    \item We establish $O(1/T)$ last-iterate convergence for exact population proximal updates and evaluate the method against matched preference-optimization controls and reward-model-based multi-objective baselines under independent held-out metrics.
\end{itemize}

\section{Background}
\label{sec:background}
\paragraph{General preference models.}
We use contextual-bandit notation. A prompt (context) $x$ is drawn from a distribution $\D$, a response $y$ is an action from a finite set $\Y_x$, and $\pi(y\mid x)$ is the policy. The reference policy is denoted by $\piref$. For the finite-dimensional population analysis, the support of $\D$ and every $\Y_x$ are finite, and
\[
\Pi=\prod_{x\in\operatorname{supp}(\D)}\Delta(\Y_x)
\]
is the contextual policy simplex. Objective $k\in[K]$ is represented by a pairwise preference model $P_k(y\succ z\mid x)\in[0,1]$ satisfying
\begin{equation}
P_k(y\succ z\mid x)+P_k(z\succ y\mid x)=1,
\qquad
P_k(y\succ y\mid x)=\tfrac12.
\label{eq:skew}
\end{equation}
We center each comparison at a tie,
\begin{equation}
\Delta_k(y,z\mid x)
=
P_k(y\succ z\mid x)-\tfrac12,
\qquad
\Delta_k(y,z\mid x)=-\Delta_k(z,y\mid x),
\label{eq:centered}
\end{equation}
and define the expected objective-$k$ payoff of learner policy $\pi$ against comparator policy $\nu$ as
\begin{equation}
J_k(\pi,\nu)
=
\E_{x\sim\D,\,y\sim\pi(\cdot\mid x),\,z\sim\nu(\cdot\mid x)}
\Delta_k(y,z\mid x).
\label{eq:game-payoff}
\end{equation}
Thus $J_k(\pi,\nu)>0$ means that $\pi$ is preferred to $\nu$ on average under objective $k$. This representation does not assume a scalar reward or transitive preferences.

\paragraph{Bargaining problems.}
A $K$-dimensional bargaining problem consists of a feasible value set $\G\subseteq\R^K$ and a disagreement point $d\in\G$, the outcome obtained if no agreement is reached. For a feasible value vector $v$, define its objective-wise surplus over disagreement as
\begin{equation}
s_k(v):=v_k-d_k,
\qquad
s(v)=(s_1(v),\ldots,s_K(v)).
\label{eq:classical-surplus}
\end{equation}
A value vector is individually rational when $s(v)\ge0$ componentwise, and strictly individually rational when $s(v)>0$. On a compact, convex, and downward-closed feasible set with at least one strictly individually rational point, the classical Nash solution is
\begin{equation}
v^{\mathrm N}
\in
\arg\max_{v\in\G:\,s(v)>0}
\sum_{k=1}^K\log s_k(v).
\label{eq:classical-nash}
\end{equation}
Equivalently, it maximizes the product of the objective-wise surpluses \citep{nash1950}.

\section{Objective-Wise Preference Games}
\label{sec:game-utility}
We first convert each objective's pairwise preference model into a scalar policy value without imposing a scalar reward on individual responses. For objective $k$, an adaptive comparator $\nu$ searches for responses against which the learner performs poorly. The reference policy $\piref$ anchors this comparator and later serves as the fallback in the bargaining problem.

\begin{definition}[Entropy-regularized game value]
\label{def:game-value}
For objective $k$ and temperature $\beta_k>0$, define
\begin{equation}
V_{k,\beta_k}(\pi)
=
\min_{\nu}
\left\{
J_k(\pi,\nu)
+
\beta_k\E_{x\sim\D}
\KL\!\left(\nu(\cdot\mid x)\,\|\,\piref(\cdot\mid x)\right)
\right\}.
\label{eq:game-value}
\end{equation}
\end{definition}
The comparator minimizes the learner's pairwise payoff, so it emphasizes responses that expose weaknesses of the current policy. The KL term prevents it from moving arbitrarily far from $\piref$. Accordingly, a large $\beta_k$ produces a local evaluation near the reference, whereas $\beta_k\downarrow0$ approaches the unregularized worst-comparator value.

For a candidate comparator response $z$, let
\begin{equation}
m_k^\pi(z\mid x)
:=
\E_{y\sim\pi(\cdot\mid x)}\Delta_k(y,z\mid x)
\label{eq:r-def}
\end{equation}
be the learner's expected margin against $z$. Smaller values identify stronger challenges. The inner problem in \eqref{eq:game-value} has the closed-form solution
\begin{align}
\nu^\star_{k,\pi}(z\mid x)
&=
\frac{
\piref(z\mid x)\exp\{-m_k^\pi(z\mid x)/\beta_k\}
}{
\E_{z'\sim\piref(\cdot\mid x)}
\exp\{-m_k^\pi(z'\mid x)/\beta_k\}
},
\label{eq:regularized-opponent}\\
V_{k,\beta_k}(\pi)
&=
-\beta_k\E_{x\sim\D}
\log
\E_{z\sim\piref(\cdot\mid x)}
\exp\{-m_k^\pi(z\mid x)/\beta_k\}.
\label{eq:soft-value}
\end{align}
Thus the opponent is obtained by reweighting reference responses toward those with smaller margins, and the game value is a soft minimum of these margins.

\begin{proposition}[Game-value structure]
\label{prop:game-structure}
For every $k$ and $\beta_k\ge0$, $V_{k,\beta_k}$ is continuous and concave in $\pi$. For $\beta_k>0$, its policy-gradient coordinate is
\begin{equation}
\left[\nabla_\pi V_{k,\beta_k}(\pi)\right]_{x,y}
=
\D(x)\,Q_k^\pi(y\mid x),
\qquad
Q_k^\pi(y\mid x)
:=
\E_{z\sim\nu^\star_{k,\pi}(\cdot\mid x)}
\Delta_k(y,z\mid x).
\label{eq:game-gradient}
\end{equation}
\end{proposition}
The proof is in Appendix~\ref{app:proof-game-structure}. The response-level score $Q_k^\pi(y\mid x)$ is the expected margin of response $y$ against the current objective-$k$ opponent; it therefore measures the local effect of increasing $\pi(y\mid x)$. An adaptive opponent can distinguish policies that have the same win rate against one fixed reference but different vulnerabilities to other responses. Appendix~\ref{app:anchor-limit} gives a finite cyclic example at the unregularized endpoint $\beta=0$; this illustration is not used in the later assumptions or proofs.

\section{Nash Bargaining over Game-Value Improvements}
\label{sec:bargaining}
We now combine the objective-wise game values using the reference policy as the common fallback. For each objective,
\begin{equation}
d_k:=V_{k,\beta_k}(\piref),
\qquad
s_k(\pi):=V_{k,\beta_k}(\pi)-d_k.
\label{eq:game-surplus}
\end{equation}
The surplus $s_k(\pi)$ is the improvement of $\pi$ over the deployed reference under objective $k$.

\begin{assumption}[Joint improvement and support]
\label{ass:joint}
The reference $\piref$ has full support, and there exists a policy $\bar\pi$ such that $s_k(\bar\pi)>0$ for every $k$.
\end{assumption}
The assumption states that at least one policy improves all objectives over the fallback. Under this condition, the NBPO population target is
\begin{equation}
\pi^{\mathrm N}
\in
\arg\max_{\pi\in\Pi_+}F(\pi),
\qquad
F(\pi):=\sum_{k=1}^K\log s_k(\pi),
\qquad
\Pi_+:=\{\pi\in\Pi:s_k(\pi)>0,\ \forall k\}.
\label{eq:nash-welfare}
\end{equation}
Equivalently, introducing auxiliary lower bounds $u_k>0$ yields
\begin{equation}
(\pi^{\mathrm N},u^{\mathrm N})
\in
\arg\max_{\pi,\,u\in\R_{++}^K}
\sum_{k=1}^K\log u_k
\quad
\text{s.t.}\quad
u_k\le s_k(\pi),\ k\in[K],
\qquad
\pi\in\Pi.
\label{eq:nash-bargaining}
\end{equation}
Because the objective is increasing in every $u_k$, each constraint is tight at an optimum and $u_k^{\mathrm N}=s_k(\pi^{\mathrm N})$.

To connect \eqref{eq:nash-welfare} to the classical bargaining problem, define
\[
V_{\boldsymbol\beta}(\pi)
=(V_{1,\beta_1}(\pi),\ldots,V_{K,\beta_K}(\pi)),
\qquad
d=(d_1,\ldots,d_K),
\]
and the downward-closed feasible value set
\begin{equation}
\G_{\boldsymbol\beta}
:=
\left\{
v\in\R^K:
d\le v\le V_{\boldsymbol\beta}(\pi)
\text{ componentwise for some }\pi\in\Pi
\right\}.
\label{eq:comprehensive-hull}
\end{equation}

\begin{proposition}[Policy-induced Nash bargaining]
\label{prop:nash-outcome}
Under Assumption~\ref{ass:joint}, $\G_{\boldsymbol\beta}$ is compact, convex, downward closed above $d$, and nondegenerate. The policies solving \eqref{eq:nash-welfare} realize the classical Nash solution of $(\G_{\boldsymbol\beta},d)$. Consequently, every optimizer has positive surplus in every objective and is Pareto optimal in game-value space; all optimizers induce the same game-value vector; and the optimizer set is unchanged by independent transformations
\[
(V_{k,\beta_k},d_k)
\mapsto
(c_kV_{k,\beta_k}+b_k,\ c_kd_k+b_k),
\qquad c_k>0.
\]
\end{proposition}
Appendix~\ref{app:nash-properties} verifies the feasible-set conditions. Once those conditions hold, the stated outcome properties are standard consequences of the classical Nash bargaining solution \citep{nash1950}; we do not claim them as new bargaining axioms. In LLM alignment, they mean that the population solution improves every objective over the reference, cannot be Pareto dominated by another feasible policy, and does not change merely because different judges report their values in different numerical units. Appendix~\ref{app:rule-witnesses} compares these guarantees with utilitarian and maxmin aggregation.

\section{Nash Bargaining Preference Optimization}
\label{sec:algorithm}
Optimizing \eqref{eq:nash-welfare} presents three practical difficulties. First, the reference has $s_k(\piref)=0$, so the logarithmic objective and its gradient are undefined there. Second, each objective's opponent depends on the policy being optimized. Third, the ideal exponential policy update contains a language-model normalizer that cannot be computed. This section addresses the three issues with a constrained proximal step, alternating opponent and policy updates, and pairwise log-ratio regression.

\subsection{Constrained proximal update and adaptive objective weights}
For $\pi\in\Pi_+$, Proposition~\ref{prop:game-structure} and the chain rule give
\begin{equation}
\left[\nabla_\pi F(\pi)\right]_{x,y}
=
\D(x)\sum_{k=1}^K
\frac{Q_k^\pi(y\mid x)}{s_k(\pi)}.
\label{eq:nash-gradient}
\end{equation}
The inverse-surplus factors in \eqref{eq:nash-gradient} emphasize objectives that have improved less. Because this expression cannot be evaluated at $\piref$, outer stage $t$ instead solves the constrained KL-proximal problem
\begin{equation}
\begin{aligned}
(\pi_{t+1},u_{t+1})\in
\arg\max_{\pi,\,u\in\R_{++}^K}
\quad&
\sum_{k=1}^K\log u_k
-\frac{1}{\eta_t}D(\pi\|\pi_t)\\
\text{s.t.}\quad&
u_k\le s_k(\pi),\qquad k\in[K],\\
&\pi\in\Pi,
\end{aligned}
\label{eq:prox-subproblem}
\end{equation}
where $\pi_0=\piref$, $\eta_t>0$, and
$D(\pi\|\pi_t)=\E_{x\sim\D}\KL(\pi(\cdot\mid x)\|\pi_t(\cdot\mid x))$.
Assumption~\ref{ass:joint} supplies a feasible positive-surplus policy, while $\piref$ serves only as the initial proximal center.

The dual of \eqref{eq:prox-subproblem} is
\begin{equation}
\min_{\lambda\in\R_{++}^K}
\phi_t(\lambda),
\qquad
\phi_t(\lambda)
=
-\sum_{k=1}^K\log\lambda_k-K
+
\max_{\pi\in\Pi}
\left\{
\sum_{k=1}^K\lambda_k s_k(\pi)
-\frac{1}{\eta_t}D(\pi\|\pi_t)
\right\}.
\label{eq:prox-dual}
\end{equation}
For fixed $\lambda$, denote the unique inner maximizer by
\begin{equation}
\pi_t^\star(\lambda)
:=
\arg\max_{\pi\in\Pi}
\left\{
\sum_{k=1}^K\lambda_k s_k(\pi)
-\frac{1}{\eta_t}D(\pi\|\pi_t)
\right\}.
\label{eq:weighted-policy}
\end{equation}
Danskin's theorem yields the low-dimensional dual gradient
\begin{equation}
\nabla_k\phi_t(\lambda)
=
s_k\!\left(\pi_t^\star(\lambda)\right)-\frac{1}{\lambda_k}.
\label{eq:dual-gradient}
\end{equation}
The complete dual derivation is in Appendix~\ref{app:prox-dual}.

\begin{proposition}[Inverse-surplus objective weights]
\label{prop:prox-dual}
At an optimal primal--dual solution of \eqref{eq:prox-subproblem},
\begin{equation}
u_{t+1,k}=s_k(\pi_{t+1}),
\qquad
\lambda_{t+1,k}
=\frac{1}{u_{t+1,k}}
=\frac{1}{s_k(\pi_{t+1})}.
\label{eq:kkt-weights}
\end{equation}
\end{proposition}
Thus the constrained problem recovers the same inverse-surplus weighting visible in \eqref{eq:nash-gradient}: objectives with smaller achieved improvements receive larger weights, without choosing those weights in advance.

\subsection{Sampled policy fitting and dual update}
Fix outer stage $t$ and a dual vector $\lambda$. Because the opponent $\nu^\star_{k,\pi}$ depends on the learner, we approximately solve \eqref{eq:weighted-policy} by alternating two operations. At inner iteration $r$, first construct $\nu^\star_{k,\pi^{(r)}}$ for each objective. Then hold these opponents fixed and fit the next policy. Define
\[
Q_k^{(r)}(y\mid x)
:=
\E_{z\sim\nu^\star_{k,\pi^{(r)}}(\cdot\mid x)}
\Delta_k(y,z\mid x).
\]
The ideal fixed-opponent update is
\begin{equation}
\pi^{(r+1)}(y\mid x)
\propto
\pi_t(y\mid x)
\exp\!\left\{
\eta_t\sum_{k=1}^K\lambda_kQ_k^{(r)}(y\mid x)
\right\}.
\label{eq:fixed-opponent-update}
\end{equation}
Its normalizer disappears when comparing two responses. Let
\begin{equation}
h_t(\pi;x,y,y')
:=
\log\frac{\pi(y\mid x)}{\pi_t(y\mid x)}
-
\log\frac{\pi(y'\mid x)}{\pi_t(y'\mid x)}.
\label{eq:log-ratio-change}
\end{equation}
Then \eqref{eq:fixed-opponent-update} implies
\begin{equation}
h_t(\pi^{(r+1)};x,y,y')
=
\eta_t\sum_{k=1}^K\lambda_k
\left(Q_k^{(r)}(y\mid x)-Q_k^{(r)}(y'\mid x)\right).
\label{eq:fixed-opponent-pairwise}
\end{equation}

To estimate the right-hand side, sample $(y,y')$ from a full-support pair proposal $\xi(\cdot,\cdot\mid x)$ and sample $z_k\sim\nu^\star_{k,\pi^{(r)}}(\cdot\mid x)$. Query independent binary judgments
\[
B_k\sim\operatorname{Bernoulli}(P_k(y\succ z_k\mid x)),
\qquad
B'_k\sim\operatorname{Bernoulli}(P_k(y'\succ z_k\mid x)),
\]
and set
\begin{equation}
Z_k^{(r)}:=B_k-B'_k.
\label{eq:binary-target}
\end{equation}
Because
\begin{equation}
\E[Z_k^{(r)}\mid x,y,y']
=
Q_k^{(r)}(y\mid x)-Q_k^{(r)}(y'\mid x),
\label{eq:binary-target-mean}
\end{equation}
the weighted binary label is an unbiased target for the desired log-ratio. We therefore fit
\begin{equation}
\mathcal L_{\mathrm{reg}}^{(t,\lambda,r)}(\pi)
=
\E\!\left[
 h_t(\pi;x,y,y')
-
\eta_t\sum_{k=1}^K\lambda_k Z_k^{(r)}
\right]^2,
\label{eq:regression-loss}
\end{equation}
where the expectation follows the sampling process above. Recomputing the opponents after each policy fit produces a fixed point that satisfies the optimality condition of \eqref{eq:weighted-policy}; Appendix~\ref{app:prox-dual} gives the population argument.

After approximately solving the weighted policy problem, we estimate its surplus vector and update the positive dual weights by
\begin{equation}
\lambda^{(m+1)}
=
\Pi_{\Lambda}\!\left[
\lambda^{(m)}
-
\gamma_m\left(
\widehat s\!\left(\pi_t^\star(\lambda^{(m)})\right)
-
(\lambda^{(m)})^{-1}
\right)
\right],
\label{eq:dual-update}
\end{equation}
where the inverse is elementwise and $\Pi_\Lambda$ keeps the numerical weights in a positive bounded box. The indices have distinct roles: $r$ alternates opponent construction and policy fitting for fixed $\lambda$; $m$ updates the dual weights within one proximal stage; and $t$ commits the next outer policy.

\begin{algorithm}[t]
\caption{Nash Bargaining Preference Optimization}
\label{alg:nbpo}
\begin{algorithmic}[1]
\REQUIRE Reference $\piref$, preference models $P_{1:K}$, temperatures $\beta_{1:K}$, outer stages $T$, dual steps $M$, policy--opponent steps $R$.
\STATE Estimate $d_k=V_{k,\beta_k}(\piref)$ and set $\pi_0\leftarrow\piref$.
\FOR{$t=0,\ldots,T-1$}
    \STATE Initialize positive weights $\lambda^{(0)}$ (warm-start from the previous stage when available).
    \FOR{$m=0,\ldots,M-1$}
        \STATE $\widehat\pi^{(m)}\leftarrow\textsc{PolicyFit}(\pi_t,\lambda^{(m)},R)$ using opponent reweighting and \eqref{eq:regression-loss}.
        \STATE Estimate $\widehat s(\widehat\pi^{(m)})$ and update $\lambda^{(m+1)}$ by \eqref{eq:dual-update}.
    \ENDFOR
    \STATE $\widetilde\pi_{t+1}\leftarrow\textsc{PolicyFit}(\pi_t,\lambda^{(M)},R)$.
    \STATE Set $\pi_{t+1}\leftarrow\widetilde\pi_{t+1}$ if all held-out empirical surpluses are positive; otherwise retain $\pi_t$.
\ENDFOR
\STATE \textbf{return} $\pi_T$.
\end{algorithmic}
\end{algorithm}
The held-out check is an implementation safeguard rather than a finite-sample guarantee. The convergence result below concerns exact population solutions of \eqref{eq:prox-subproblem}.

\subsection{Last-iterate convergence}
Because Algorithm~\ref{alg:nbpo} returns the final outer policy rather than an average of past policies, we analyze the last population iterate.
\begin{theorem}[Last-iterate convergence]
\label{thm:last-iterate}
Under Assumption~\ref{ass:joint}, let $\pi^\star$ maximize \eqref{eq:nash-welfare}, set $\pi_0=\piref$, and solve \eqref{eq:prox-subproblem} exactly at every outer iteration with arbitrary $\eta_t>0$. Then $F(\pi_{t+1})\ge F(\pi_t)$ for $t\ge1$, and for every $T\ge1$,
\begin{equation}
F(\pi^\star)-F(\pi_T)
\le
\frac{D(\pi^\star\|\piref)}{\sum_{t=0}^{T-1}\eta_t}.
\label{eq:last-iterate-rate}
\end{equation}
\end{theorem}
A constant proximal step therefore gives an $O(1/T)$ final-iterate Nash-welfare gap. Appendix~\ref{app:convergence} gives the proof.

\FloatBarrier
\section{Experiments}
\label{sec:experiments}
\paragraph{Setup.}
We evaluate fixed objective panels on \textsc{UltraFeedback}, \textsc{TL;DR}, and \textsc{SafeRLHF}. The primary backbone and reference policy are \texttt{Llama-3.1-8B-Instruct}. A prompted \texttt{Qwen3-32B} model supplies objective-specific training comparisons, queried in both presentation orders, while \texttt{Llama-3.3-70B-Instruct} and \texttt{Phi-4} provide independent held-out evaluation. All primary comparisons use prompt-disjoint train, validation, and test splits and report the mean and standard deviation over three end-to-end training seeds. Appendix~\ref{app:hyperparameters} records data, generation, optimization, model-selection, and compute details; Appendix~\ref{app:labeling} specifies the preference-labeling interface.

\paragraph{Baselines.}
We organize the comparison around the two design choices in \NBPO{}. Three Nash-aggregated controls change only the objective representation: adaptive preference-game values (\NBPO{}), fixed-reference pairwise margins (Fixed-reference Nash), and objective-specific Bradley--Terry rewards (BT-RM--Nash). Two controls retain the adaptive game representation but replace Nash aggregation with equal-weight utilitarian aggregation or a Kalai--Smorodinsky proportional-gains rule (Game-utilitarian and Game-KS) \citep{kalai1975}. BT-RM--utilitarian completes the matched representation--aggregation comparison. External multi-objective baselines are PROSPER/MaxEntBW, MOPO, and RACO \citep{zhang2026prosper,agnihotri2026mopo,chen2026raco}. Every method uses the same prompt splits, response pool, comparison bank, policy-training budget, and held-out evaluator whenever its algorithm permits; family-valued methods select one operating point using validation data only.

\paragraph{Evaluation.}
Our primary cross-method metric is the independent objective-wise win rate against the common reference,
\begin{equation}
W_k^{\mathrm{eval}}(\pi)
:=
\E_{x\sim\D,\,y\sim\pi(\cdot\mid x),\,y_{\mathrm{ref}}\sim\piref(\cdot\mid x)}
P_k^{\mathrm{eval}}(y\succ y_{\mathrm{ref}}\mid x).
\label{eq:external-win-rate}
\end{equation}
All methods are ranked by the same evaluators, so this metric does not privilege the \NBPO{} training objective. Table~\ref{tab:general-benchmark-completion} reports general capability for the \textsc{UltraFeedback}-trained checkpoints. Table~\ref{tab:multi-dataset-external} summarizes the worst and average objective-wise win rates on all three panels, and Appendix~\ref{app:objectivewise} lists every objective. Table~\ref{tab:component-ablation} isolates objective representation from aggregation on a shared \textsc{UltraFeedback} protocol.

% CLAUDE TABLE CONTRACT: when present, the generated file must contain the
% complete table* environment and preserve label tab:general-benchmark-completion.
\IfFileExists{\NBPOResultDir/table1_generated.tex}{%
  % Generated by scripts/experiments/iclr2027_table1_v2/build_results_tables.py
% Do not edit by hand: regenerate from the run registry.
% \pending{} marks a cell whose run has not completed all three seeds.
% \providecommand{\pending}{\textit{pending}}
% rows complete: 0/10

\begin{tabular}{lcc}
\toprule
Method & Worst objective & Average objective \\
\midrule
Base ($\pi_\mathrm{ref}$) & \pending{} & \pending{} \\
NBPO & \pending{} & \pending{} \\
Fixed-reference Nash & \pending{} & \pending{} \\
BT-RM-Nash & \pending{} & \pending{} \\
Game-utilitarian & \pending{} & \pending{} \\
Game-KS & \pending{} & \pending{} \\
BT-RM-utilitarian & \pending{} & \pending{} \\
PROSPER / MaxEntBW & \pending{} & \pending{} \\
MOPO & \pending{} & \pending{} \\
RACO & \pending{} & \pending{} \\
\bottomrule
\end{tabular}
%
}{%
\begin{table*}[t]
\centering\scriptsize
\setlength{\tabcolsep}{3.2pt}
\caption{General-capability evaluation of \textsc{UltraFeedback}-trained \texttt{Llama-3.1-8B-Instruct} policies. Pairwise columns are held-out win rates against the common reference; IFEval and TruthfulQA use their official metrics. Entries are mean $\pm$ standard deviation over three training seeds. Higher is better.}
\label{tab:general-benchmark-completion}
\begin{tabular}{lccccc}
\toprule
method & Arena-Hard v2 ref-win & AlpacaEval 2 LC & MT-Bench ref-win & IFEval strict & TruthfulQA MC2 \\
\midrule
Base $=\piref=\pi_0$ & \pending & \pending & \pending & \pending & \pending \\
\midrule
\NBPO{} (ours) & \pending & \pending & \pending & \pending & \pending \\
Fixed-reference Nash & \pending & \pending & \pending & \pending & \pending \\
BT-RM--Nash & \pending & \pending & \pending & \pending & \pending \\
Game-utilitarian & \pending & \pending & \pending & \pending & \pending \\
Game-KS & \pending & \pending & \pending & \pending & \pending \\
BT-RM--utilitarian & \pending & \pending & \pending & \pending & \pending \\
\midrule
PROSPER/MaxEntBW & \pending & \pending & \pending & \pending & \pending \\
MOPO & \pending & \pending & \pending & \pending & \pending \\
RACO & \pending & \pending & \pending & \pending & \pending \\
\bottomrule
\end{tabular}
\end{table*}%
}

\paragraph{Primary multi-objective comparison.}
Table~\ref{tab:multi-dataset-external} reports the worst and average $W_k^{\mathrm{eval}}$ in each panel. The row order is fixed before evaluation, and incomplete cells remain unreported rather than being backfilled from earlier checkpoints or test-selected operating points.

% CLAUDE TABLE CONTRACT: when present, the generated file must contain the
% complete table* environment and preserve label tab:multi-dataset-external.
\IfFileExists{\NBPOResultDir/table2_generated.tex}{%
  % Generated by scripts/experiments/iclr2027_table1_v2/build_results_tables.py
% Do not edit by hand: regenerate from the run registry.
% \pending{} marks a cell whose run has not completed all three seeds.
% \providecommand{\pending}{\textit{pending}}
% rows complete: 0/10

\begin{tabular}{lcccc}
\toprule
Method & instruction following & truthfulness & honesty & helpfulness \\
\midrule
Base ($\pi_\mathrm{ref}$) & \pending{} & \pending{} & \pending{} & \pending{} \\
NBPO & \pending{} & \pending{} & \pending{} & \pending{} \\
Fixed-reference Nash & \pending{} & \pending{} & \pending{} & \pending{} \\
BT-RM-Nash & \pending{} & \pending{} & \pending{} & \pending{} \\
Game-utilitarian & \pending{} & \pending{} & \pending{} & \pending{} \\
Game-KS & \pending{} & \pending{} & \pending{} & \pending{} \\
BT-RM-utilitarian & \pending{} & \pending{} & \pending{} & \pending{} \\
PROSPER / MaxEntBW & \pending{} & \pending{} & \pending{} & \pending{} \\
MOPO & \pending{} & \pending{} & \pending{} & \pending{} \\
RACO & \pending{} & \pending{} & \pending{} & \pending{} \\
\bottomrule
\end{tabular}
%
}{%
\begin{table*}[t]
\centering\scriptsize
\setlength{\tabcolsep}{2.5pt}
\caption{Primary multi-objective evaluation under independent objective-wise win rates against the common reference. We report the worst and average objective win rate in each panel. Entries are mean $\pm$ standard deviation over three training seeds. Higher is better.}
\label{tab:multi-dataset-external}
\begin{tabular}{lcccccc}
\toprule
& \multicolumn{2}{c}{\textsc{UltraFeedback}} & \multicolumn{2}{c}{\textsc{TL;DR}} & \multicolumn{2}{c}{\textsc{SafeRLHF}}\\
\cmidrule(lr){2-3}\cmidrule(lr){4-5}\cmidrule(lr){6-7}
method & worst $W^{\rm eval}$ & avg. $W^{\rm eval}$ & worst $W^{\rm eval}$ & avg. $W^{\rm eval}$ & worst $W^{\rm eval}$ & avg. $W^{\rm eval}$\\
\midrule
Base $=\piref$ & \pending & \pending & \pending & \pending & \pending & \pending \\
\midrule
\NBPO{} (ours) & \pending & \pending & \pending & \pending & \pending & \pending \\
Fixed-reference Nash & \pending & \pending & \pending & \pending & \pending & \pending \\
BT-RM--Nash & \pending & \pending & \pending & \pending & \pending & \pending \\
Game-utilitarian & \pending & \pending & \pending & \pending & \pending & \pending \\
Game-KS & \pending & \pending & \pending & \pending & \pending & \pending \\
BT-RM--utilitarian & \pending & \pending & \pending & \pending & \pending & \pending \\
\midrule
PROSPER/MaxEntBW & \pending & \pending & \pending & \pending & \pending & \pending \\
MOPO & \pending & \pending & \pending & \pending & \pending & \pending \\
RACO & \pending & \pending & \pending & \pending & \pending & \pending \\
\bottomrule
\end{tabular}
\end{table*}%
}

\paragraph{Separating representation and aggregation.}
Table~\ref{tab:component-ablation} changes one design choice at a time. \NBPO{} versus Fixed-reference Nash and BT-RM--Nash holds Nash aggregation fixed while changing the within-objective representation. \NBPO{} versus Game-utilitarian and Game-KS holds the adaptive game representation fixed while changing the cross-objective compromise rule. BT-RM--Nash versus BT-RM--utilitarian repeats the aggregation comparison under scalar rewards.

% CLAUDE TABLE CONTRACT: when present, the generated file must contain the
% complete table* environment and preserve label tab:component-ablation.
\IfFileExists{\NBPOResultDir/table3_generated.tex}{%
  % Generated by scripts/experiments/iclr2027_table1_v2/build_results_tables.py
% Do not edit by hand: regenerate from the run registry.
% \pending{} marks a cell whose run has not completed all three seeds.
% \providecommand{\pending}{\textit{pending}}
% rows complete: 0/10

\begin{tabular}{lccc}
\toprule
Matched comparison & $\Delta$ worst & 95\% CI & Holm $p$ \\
\midrule
NBPO vs Fixed-reference Nash & \pending{} & \pending{} & \pending{} \\
NBPO vs BT-RM-Nash & \pending{} & \pending{} & \pending{} \\
NBPO vs Game-utilitarian & \pending{} & \pending{} & \pending{} \\
NBPO vs Game-KS & \pending{} & \pending{} & \pending{} \\
NBPO vs PROSPER / MaxEntBW & \pending{} & \pending{} & \pending{} \\
BT-RM-Nash vs BT-RM-utilitarian & \pending{} & \pending{} & \pending{} \\
\bottomrule
\end{tabular}
%
}{%
\begin{table*}[t]
\centering\footnotesize
\setlength{\tabcolsep}{4.0pt}
\caption{Matched component comparison on the same held-out \textsc{UltraFeedback} prompts. All variants share the response pool, comparison bank, policy-training budget, and evaluator. Entries are mean $\pm$ standard deviation over three training seeds. Higher is better.}
\label{tab:component-ablation}
\begin{tabular}{lllcc}
\toprule
method & objective representation & aggregation & worst $W^{\rm eval}$ & avg. $W^{\rm eval}$\\
\midrule
\NBPO{} (ours) & adaptive preference game & Nash & \pending & \pending \\
Fixed-reference Nash & fixed-reference margin & Nash & \pending & \pending \\
BT-RM--Nash & scalar BT reward & Nash & \pending & \pending \\
Game-utilitarian & adaptive preference game & utilitarian & \pending & \pending \\
Game-KS & adaptive preference game & Kalai--Smorodinsky & \pending & \pending \\
BT-RM--utilitarian & scalar BT reward & utilitarian & \pending & \pending \\
\bottomrule
\end{tabular}
\end{table*}%
}

\paragraph{Diagnostics.}
The adaptive game-value surpluses $s_k$ are method-aligned diagnostics rather than the headline cross-method metric. Appendix~\ref{app:optimization-diagnostics} therefore reports finite-pool realization, regression fit, scale-rescaling behavior, and cycle-conditioned comparisons separately from Tables~\ref{tab:general-benchmark-completion}--\ref{tab:component-ablation}.

\section{Related Work}
\label{sec:related}
\paragraph{General pairwise preferences.}
NLHF and its scalable variants optimize directly against general pairwise preference models rather than fitting scalar rewards \citep{munos2024nash,wu2025sppo,zhang2025iterative,calandriello2024online,rosset2024direct,zhou2025egpo}. Blackwell-style methods extend this view to vector-valued feedback and target sets \citep{blackwell1956analog,bhatia2020multicriteria}. PROSPER introduces a maximum-entropy Blackwell winner with a scalable single-policy update and an inner worst-case choice over objective weights and opponents \citep{zhang2026prosper}, while RACO combines objective-specific direct-preference losses through conflict-aware gradients \citep{chen2026raco}. \NBPO{} instead keeps one regularized game value per objective and applies Nash bargaining across the improvements over a shared fallback.

\paragraph{Reward-based multi-objective alignment.}
MODPO, Rewarded Soups, and Panacea expose user-specified trade-offs through scalarization, policy interpolation, or preference conditioning \citep{zhou2024beyond,rame2023rewarded,zhong2024panacea}. MaxMin-RLHF and RMOD optimize weakest-objective criteria, while MOPO designates a primary objective and lower-bound constraints for the others \citep{chakraborty2024maxmin,son2026rmod,agnihotri2026mopo}. Our matched BT-RM controls instead use the same objective-specific comparison bank and policy-fitting budget as \NBPO{}, allowing scalar reward representation to be changed without simultaneously changing the aggregation rule or deployment target.

\paragraph{Social choice and bargaining.}
The Nash solution and its Pareto, uniqueness, and affine-invariance properties are classical \citep{nash1950}. The Kalai--Smorodinsky solution instead equalizes normalized gains toward an ideal point and provides a strong proportional-gains comparison \citep{kalai1975}. Social-choice work argues that alignment should state its aggregation rule explicitly \citep{conitzer2024social}. Multi-party RLHF studies utilitarian, Nash, and leximin welfare over party-specific learned rewards and separately considers reward-free von Neumann solutions \citep{zhong2024nashwelfare}; Nash bargaining has also been used to combine task gradients in multi-task learning \citep{navon2022nashmtl}. Our contribution is not a new bargaining axiom, but a formulation and practical optimizer for applying the classical Nash rule to policy-dependent objective-wise preference-game values.

\section{Conclusion and Limitations}
\label{sec:conclusion}
\NBPO{} separates two decisions in multi-objective preference alignment: representing possibly cyclic comparisons within each objective and selecting one compromise across objectives. It uses regularized preference games for the first decision and the classical Nash bargaining rule for the second. The bargaining guarantees are inherited once the policy-induced game-value set satisfies the standard conditions; our main contribution is the constrained and partition-free optimization procedure that makes this construction trainable from binary comparisons. The exact population proximal iteration achieves an $O(1/T)$ last-iterate Nash-welfare gap.

Several limitations remain. Whether adaptive game values improve over fixed-reference or scalar-reward representations depends on the prevalence of intransitivity and on how accurately the finite-pool target is realized by the neural policy; we therefore report cycle-conditioned and regression-realization diagnostics in addition to aggregate win rates. The neural algorithm still approximates the policy--opponent fixed point, dual solution, and regression target from finite samples, whereas the convergence theorem assumes exact population updates. The joint-improvement assumption can also fail when no policy improves every objective over the chosen reference; in that case a positive-surplus Nash solution does not exist. Improvements in game value should not be interpreted as per-prompt safety guarantees.

\subsection*{AI use statement}
In this work, generative AI tools were used to critique the problem formulation, assist with mathematical derivations and proof presentation, propose experimental checks, and edit manuscript text. The authors are responsible for independently verifying every theorem, citation, implementation, empirical result, and disclosure in the final submission.

\subsection*{Ethics statement}
This work studies aggregation rules for language-model alignment, including harmlessness as one objective. Improvement in a game value is not an absolute safety certificate. Safety-critical deployment should retain hard constraints and independent safety evaluations in addition to the bargaining objective.

\subsection*{Reproducibility statement}
Algorithm~\ref{alg:nbpo} specifies the constrained training procedure. Appendix~\ref{app:proofs} provides the proximal-dual derivation and population convergence proof, while Appendices~\ref{app:hyperparameters}, \ref{app:optimization-diagnostics}, and~\ref{app:labeling} provide the matched experimental protocol, diagnostics, and judge details. A submission package should include anonymized code, exact prompts, random seeds, model checkpoints, opponent pools, dual traces, and the sampling budgets used for the regression and surplus estimates.

\bibliography{iclr2027_conference}
\bibliographystyle{iclr2027_conference}

\appendix

\section{Proofs}
\label{app:proofs}

\subsection{Proof of Proposition~\ref{prop:game-structure}}
\label{app:proof-game-structure}
For a fixed comparator $\nu$, define
\[
\mathcal J_{k,\nu}(\pi)
:=
J_k(\pi,\nu)
+
\beta_k\E_{x\sim\D}
\KL\!\left(\nu(\cdot\mid x)\,\|\,\piref(\cdot\mid x)\right).
\]
The first term is linear in $\pi$ and the second does not depend on $\pi$, so $\mathcal J_{k,\nu}$ is affine. The pointwise minimum of affine functions is concave; hence $V_{k,\beta_k}$ is concave. Since $|\Delta_k|\le1/2$,
\[
|V_{k,\beta_k}(\pi)-V_{k,\beta_k}(\pi')|
\le
\frac12\E_{x\sim\D}
\|\pi(\cdot\mid x)-\pi'(\cdot\mid x)\|_1,
\]
which also proves continuity.

For $\beta_k>0$, the inner minimization separates across prompts. Fix $x$ and abbreviate $m(z)=m_k^\pi(z\mid x)$. The prompt-wise problem is
\[
\min_{\nu\in\Delta(\Y_x)}
\sum_z\nu(z)m(z)
+
\beta_k\sum_z\nu(z)
\log\frac{\nu(z)}{\piref(z\mid x)}.
\]
The first-order condition under $\sum_z\nu(z)=1$ gives
\[
\nu(z)\propto\piref(z\mid x)\exp\{-m(z)/\beta_k\},
\]
which yields \eqref{eq:regularized-opponent}; substituting it back yields \eqref{eq:soft-value}. The KL term makes the minimizer unique. Danskin's theorem therefore gives, for every feasible direction $h$ satisfying $\sum_yh(y\mid x)=0$,
\[
D V_{k,\beta_k}(\pi)[h]
=
\E_{x\sim\D}\sum_y h(y\mid x)
\E_{z\sim\nu^\star_{k,\pi}(\cdot\mid x)}
\Delta_k(y,z\mid x).
\]
Thus the coordinate derivative is
\[
\left[\nabla_\pi V_{k,\beta_k}(\pi)\right]_{x,y}
=
\D(x)Q_k^\pi(y\mid x),
\]
proving \eqref{eq:game-gradient}.

\subsection{Connection to the Fixed-Reference}
\label{app:anchor-limit}
This subsection gives a finite-action illustration at the unregularized endpoint $\beta=0$. It is not used in Assumption~\ref{ass:joint}, the full-support analysis, or the finite-temperature experiments.

Define the fixed-reference improvement
\begin{equation}
a_k(\pi)=J_k(\pi,\piref)=P_k(\pi\succ\piref)-\tfrac12.
\label{eq:anchor-margin}
\end{equation}
Skew-symmetry gives $a_k(\piref)=0$.

\paragraph{Full cyclic example.}
For one prompt, let $\piref=\delta_r$, where $\delta_y$ is the deterministic
policy that always outputs $y$. Suppose $A$, $B$, and $C$ each beat $r$
with probability $0.75$ and form the deterministic cycle
\[
A\succ B,\qquad B\succ C,\qquad C\succ A.
\]
Both
\[
\pi_{\rm pure}=\delta_A
\qquad\text{and}\qquad
\pi_{\rm mix}=\frac{\delta_A+\delta_B+\delta_C}{3}
\]
have the same fixed-reference improvement,
\[
a(\pi_{\rm pure})
=
a(\pi_{\rm mix})
=
0.75-\frac12
=
0.25.
\]

Their worst-opponent values are nevertheless different. For the pure
policy, response $C$ defeats $A$ deterministically, so
\[
V_0(\pi_{\rm pure})
=
\min_{\nu} J(\pi_{\rm pure},\nu)
=
J(\delta_A,\delta_C)
=
-\frac12.
\]

For the uniform mixture, consider first the pure opponent $\delta_A$.
The mixture outputs $A$, $B$, and $C$ with equal probability. Against
$A$, these responses respectively tie, lose, and win, giving centered
payoffs $0$, $-1/2$, and $+1/2$. Hence
\[
J(\pi_{\rm mix},\delta_A)
=
\frac13
\left(
0-\frac12+\frac12
\right)
=
0.
\]
By symmetry,
\[
J(\pi_{\rm mix},\delta_B)
=
J(\pi_{\rm mix},\delta_C)
=
0.
\]
Against the reference,
\[
J(\pi_{\rm mix},\delta_r)=0.25,
\]
because every response in the mixture beats $r$ with probability $0.75$.
Since $J(\pi_{\rm mix},\nu)$ is linear in $\nu$, no mixed opponent can
achieve a value below the minimum over these pure opponents. Therefore
\[
V_0(\pi_{\rm mix})
=
\min_{\nu}J(\pi_{\rm mix},\nu)
=
\min\{0,0,0,0.25\}
=
0.
\]

Finally, the reference itself has
\[
V_0(\piref)
=
\min_{\nu}J(\delta_r,\nu)
=
-0.25,
\]
since $r$ loses to each of $A$, $B$, and $C$ with probability $0.75$.
Thus
\[
\begin{array}{c|ccc}
& a(\pi) & V_0(\pi) & V_0(\pi)-V_0(\piref)\\
\hline
\pi_{\rm pure} & 0.25 & -0.50 & -0.25\\
\pi_{\rm mix}  & 0.25 & 0     &  0.25
\end{array}
\]
The fixed reference therefore assigns the two policies the same score,
even though the pure policy is fully exploitable by one response whereas
the balanced mixture guarantees an average tie against every response
in the cycle.

\subsection{Policy-Induced Nash Bargaining and Classical Consequences}
\label{app:nash-properties}
We prove Proposition~\ref{prop:nash-outcome}. Since $\Pi$ is compact and each $V_{k,\beta_k}$ is continuous, $V_{\boldsymbol\beta}(\pi)$ is bounded. The definition of $\G_{\boldsymbol\beta}$ then implies boundedness. Closedness follows by taking any convergent sequence $v_n\in\G_{\boldsymbol\beta}$, choosing policies $\pi_n$ with $d\le v_n\le V_{\boldsymbol\beta}(\pi_n)$, extracting a convergent policy subsequence, and applying continuity. Downward closure is immediate from the definition.

For convexity, take $v^1,v^2\in\G_{\boldsymbol\beta}$ with witnesses $\pi_1,\pi_2$. Objective-wise concavity from Proposition~\ref{prop:game-structure} gives, for $\alpha\in[0,1]$,
\[
V_{\boldsymbol\beta}(\alpha\pi_1+(1-\alpha)\pi_2)
\ge
\alpha V_{\boldsymbol\beta}(\pi_1)+(1-\alpha)V_{\boldsymbol\beta}(\pi_2)
\]
componentwise, so $\alpha v^1+(1-\alpha)v^2\in\G_{\boldsymbol\beta}$. Assumption~\ref{ass:joint} provides a point strictly above $d$, proving nondegeneracy.

For any $v\in\G_{\boldsymbol\beta}$ with $v>d$, some policy $\pi$ satisfies $v\le V_{\boldsymbol\beta}(\pi)$. Since the Nash objective is coordinatewise increasing,
\[
\sum_k\log(v_k-d_k)
\le
\sum_k\log\bigl(V_{k,\beta_k}(\pi)-d_k\bigr)
=F(\pi).
\]
Conversely, every $\pi\in\Pi_+$ induces a feasible value vector $V_{\boldsymbol\beta}(\pi)$. Hence \eqref{eq:nash-welfare} realizes the classical Nash solution over $(\G_{\boldsymbol\beta},d)$.

The remaining statements are standard Nash-bargaining consequences, but we record the short arguments for completeness. Since the log objective tends to $-\infty$ when any surplus approaches zero, every optimizer lies strictly above $d$. Coordinatewise monotonicity implies Pareto optimality. Strict concavity of $\sum_k\log s_k$ over value vectors implies that all optimal policies induce the same surplus, hence the same game-value vector, although the policy itself can be nonunique.

Finally, under independent positive affine transformations
\[
V_k'(\pi)=c_kV_k(\pi)+b_k,
\qquad
d_k'=c_kd_k+b_k,
\qquad c_k>0,
\]
the transformed surplus is $s_k'(\pi)=c_ks_k(\pi)$ and
\[
\sum_k\log s_k'(\pi)
=
\sum_k\log s_k(\pi)+\sum_k\log c_k.
\]
The added term is independent of $\pi$, so the optimizer set is unchanged.

\begin{table}[htbp]
\centering\footnotesize
\setlength{\tabcolsep}{4.2pt}
\caption{Classical population guarantees of common aggregation rules under joint improvement. The Nash row follows from the standard bargaining solution; the other rows are supported by the finite counterexamples below.}
\label{tab:rule-guarantees}
\begin{tabular}{lcccc}
\toprule
rule & strict IR & Pareto optimal & unique value & affine invariant\\
\midrule
positive-weight utilitarian & no & yes & no & no\\
absolute maxmin & no & no & no & no\\
surplus maxmin & yes & no & no & no\\
Nash bargaining & yes & yes & yes & yes\\
\bottomrule
\end{tabular}
\end{table}

\subsection{Counterexamples for Alternative Aggregation Rules}
\label{app:rule-witnesses}

The positive guarantees in Table~\ref{tab:rule-guarantees} follow
directly from the corresponding objectives. Because $w_k>0$, the
positive-weight utilitarian objective $\sum_k w_k V_k(\pi)$ is strictly
increasing in every game-value coordinate, so none of its optimizers can
be Pareto dominated. Under Assumption~\ref{ass:joint}, surplus maxmin
also attains a strictly positive minimum surplus because a jointly
improving policy exists. The examples below establish the guarantees
that can fail for the alternative rules.

All examples use three responses, so a policy is a probability vector
$\pi=(\pi_1,\pi_2,\pi_3)\in\Delta_3$, and use the uniform reference
$\piref=(1/3,1/3,1/3)$. For objective $k$, $\Delta_k$ is the skew-symmetric
pairwise payoff matrix: $(\Delta_k)_{ij}$ is the centered preference margin of response
$i$ against response $j$. At $\beta_k=0$,
\[
V_k(\pi)
=
\min_j(\pi^\top \Delta_k)_j,
\qquad
d_k=V_k(\piref),
\qquad
s_k(\pi)=V_k(\pi)-d_k.
\]
Thus $\pi^\top \Delta_k$ contains the expected payoff of $\pi$ against each
pure opponent, and $V_k(\pi)$ is its smallest coordinate.

Because each $V_k$ is the minimum of linear functions, the optimizers
below are obtained from small linear programs. For example,
equal-weight utilitarian optimization can be written as
\[
\max_{\pi,t_1,t_2}\ t_1+t_2
\quad\text{s.t.}\quad
t_k\le(\pi^\top \Delta_k)_j
\ \ \forall k,j,\qquad
\pi\in\Delta_3.
\]
Absolute maxmin uses one variable $t$ with
$t\le(\pi^\top \Delta_k)_j$, while surplus maxmin uses
$t+d_k\le(\pi^\top \Delta_k)_j$ for every $k,j$.

\paragraph{Utilitarian failure of strict improvement.}
Let
\[
\Delta_1=
\begin{pmatrix}
0&-1/2&-1/2\\
1/2&0&-1/2\\
1/2&1/2&0
\end{pmatrix},
\qquad
\Delta_2=
\begin{pmatrix}
0&-1/2&1/4\\
1/2&0&-1/2\\
-1/4&1/2&0
\end{pmatrix}.
\]
For the uniform reference,
\[
\piref^\top \Delta_1=(1/3,0,-1/3),
\qquad
\piref^\top \Delta_2=(1/12,0,-1/12),
\]
and hence
\[
d=(-1/3,-1/12).
\]

First, joint improvement is feasible. The policy
\[
\bar\pi=(1/4,1/4,1/2)
\]
is only a feasibility witness; it gives
\[
\bar\pi^\top \Delta_1=(3/8,1/8,-1/4),
\qquad
\bar\pi^\top \Delta_2=(0,1/8,-1/16).
\]
Therefore
\[
V(\bar\pi)=(-1/4,-1/16),
\qquad
s(\bar\pi)
=
V(\bar\pi)-d
=
(1/12,1/48)>0.
\]

In contrast, solving the equal-weight utilitarian problem
$\max_\pi V_1(\pi)+V_2(\pi)$ gives
\[
\pi_{\rm U}=(0,1/5,4/5).
\]
For this policy,
\[
\pi_{\rm U}^\top \Delta_1=(1/2,2/5,-1/10),
\qquad
\pi_{\rm U}^\top \Delta_2=(-1/10,2/5,-1/10),
\]
so
\[
V(\pi_{\rm U})=(-1/10,-1/10),
\qquad
s(\pi_{\rm U})
=
(7/30,-1/60).
\]
Thus joint improvement is possible, yet a utilitarian optimum falls
below the fallback on objective~2.

\paragraph{Absolute-maxmin failure of strict improvement.}
Let
\[
\Delta_1=
\begin{pmatrix}
0&-1/2&1/4\\
1/2&0&1/2\\
-1/4&-1/2&0
\end{pmatrix},
\qquad
\Delta_2=
\begin{pmatrix}
0&1/4&-1/2\\
-1/4&0&1/2\\
1/2&-1/2&0
\end{pmatrix}.
\]
Again,
\[
d=(-1/3,-1/12).
\]
Joint improvement is feasible, as witnessed by
\[
\bar\pi=(1/3,7/15,1/5),
\]
for which
\[
V(\bar\pi)=(-4/15,-1/60)
\]
and therefore
\[
s(\bar\pi)=(1/15,1/15)>0.
\]

Absolute maxmin instead solves
\[
\max_{\pi}\min\{V_1(\pi),V_2(\pi)\},
\]
whose linear-program solution is
\[
\pi_{\rm AM}=(0,4/5,1/5).
\]
For this policy,
\[
\pi_{\rm AM}^\top \Delta_1=(7/20,-1/10,2/5),
\qquad
\pi_{\rm AM}^\top \Delta_2=(-1/10,-1/10,2/5),
\]
so
\[
V(\pi_{\rm AM})=(-1/10,-1/10)
\]
and
\[
s(\pi_{\rm AM})
=
(7/30,-1/60).
\]
Thus absolute maxmin maximizes the smallest \emph{absolute} game value
while allowing objective~2 to fall below its own fallback.

\paragraph{A Pareto-dominated surplus-maxmin optimum.}
Let
\[
\Delta_1=
\begin{pmatrix}
0&-1/2&1/2\\
1/2&0&1/2\\
-1/2&-1/2&0
\end{pmatrix},
\qquad
\Delta_2=
\begin{pmatrix}
0&0&1/2\\
0&0&1/2\\
-1/2&-1/2&0
\end{pmatrix}.
\]
The uniform reference gives
\[
d=(-1/3,-1/6).
\]
For objective~2, no policy can achieve $V_2(\pi)>0$, while every mixture
supported on the first two responses achieves $V_2(\pi)=0$.
Consequently,
\[
s_2(\pi)\le 1/6
\]
for every policy, so the optimal surplus-maxmin value cannot exceed
$1/6$.

The two policies
\[
\pi_c=(1/3,2/3,0),
\qquad
\pi_e=(0,1,0)
\]
satisfy
\[
V(\pi_c)=(-1/6,0),
\qquad
s(\pi_c)=(1/6,1/6),
\]
and
\[
V(\pi_e)=(0,0),
\qquad
s(\pi_e)=(1/3,1/6).
\]
Both therefore maximize
\[
\min_k s_k(\pi)
\]
at $1/6$. However,
\[
s(\pi_e)>s(\pi_c)
\]
in the first coordinate and is equal in the second, so
$\pi_e$ Pareto dominates $\pi_c$. Thus the surplus-maxmin optimizer set
can contain Pareto-dominated policies.

\subsection{Dual Derivation and Regression Fixed Point}
\label{app:prox-dual}
For multipliers $\lambda_k\ge0$, the Lagrangian of \eqref{eq:prox-subproblem} is
\[
\mathcal L_t(\pi,u,\lambda)
=
\sum_k\log u_k
+
\sum_k\lambda_k\bigl(s_k(\pi)-u_k\bigr)
-
\eta_t^{-1}D(\pi\|\pi_t).
\]
For fixed $\lambda_k>0$, maximizing $\log u_k-\lambda_ku_k$ gives $u_k=1/\lambda_k$ and value $-\log\lambda_k-1$; if $\lambda_k=0$, the maximum is unbounded. Substitution yields \eqref{eq:prox-dual}, and the remaining maximization over $\pi$ is \eqref{eq:weighted-policy}. Uniqueness of that maximizer and Danskin's theorem yield \eqref{eq:dual-gradient}.

Assumption~\ref{ass:joint} provides strict feasibility: choose $\bar\pi$ with $s_k(\bar\pi)>0$ and $0<\bar u_k<s_k(\bar\pi)$. Slater's condition therefore gives strong duality. Stationarity in $u_k$ gives
\[
\lambda_{t+1,k}=\frac{1}{u_{t+1,k}}>0,
\]
and complementary slackness then forces $u_{t+1,k}=s_k(\pi_{t+1})$. This proves Proposition~\ref{prop:prox-dual}.

\paragraph{Regression fixed point.}
Fix $\lambda$ and let $\pi_t^\star(\lambda)$ solve \eqref{eq:weighted-policy}. First-order optimality gives, for every prompt and response pair in the support,
\begin{equation}
h_t(\pi_t^\star(\lambda);x,y,y')
=
\eta_t\sum_k\lambda_k
\left(
Q_k^{\pi_t^\star(\lambda)}(y\mid x)
-
Q_k^{\pi_t^\star(\lambda)}(y'\mid x)
\right).
\label{eq:prox-pairwise-kkt}
\end{equation}
At a fixed point of the alternating inner solve, the opponents in \eqref{eq:regression-loss} are exactly $\nu^\star_{k,\pi_t^\star(\lambda)}$. Equation~\eqref{eq:binary-target-mean} therefore makes the conditional mean of the binary target equal to the right-hand side of \eqref{eq:prox-pairwise-kkt}. The squared-loss conditional-mean property implies that a population minimizer of \eqref{eq:regression-loss} satisfies this pairwise identity. Full support of the pair proposal identifies all log-probability differences; normalization then recovers $\pi_t^\star(\lambda)$.

\subsection{Proof of Theorem~\ref{thm:last-iterate}}
\label{app:convergence}
Let $h(\pi)=\E_{x\sim\D}\sum_y\pi(y\mid x)\log\pi(y\mid x)$, so its Bregman divergence is $D(\pi\|\pi_t)$. Because the objective is increasing in $u$, the policy part of \eqref{eq:prox-subproblem} is equivalently
\begin{equation}
\pi_{t+1}
\in
\arg\max_{\pi\in\Pi_+}
\left\{F(\pi)-\frac1{\eta_t}D(\pi\|\pi_t)\right\}.
\label{eq:implicit-prox}
\end{equation}
The feasible set is nonempty by Assumption~\ref{ass:joint}, and $F(\pi)\to-\infty$ as any surplus approaches zero. Thus the maximizer exists in $\Pi_+$ even when the center $\pi_t$ is the reference $\piref\notin\Pi_+$. Because $\pi_t$ has full support and the remaining objective has finite directional derivatives at an interior surplus point, the inward derivative of $-D(\pi\|\pi_t)$ excludes zero-probability coordinates. Induction therefore gives a full-support, positive-surplus $\pi_{t+1}$.

Let $g_{t+1}$ be a supergradient of $F$ at $\pi_{t+1}$. First-order optimality of \eqref{eq:implicit-prox} and concavity of $F$ imply, for every $\pi\in\Pi_+$,
\begin{align}
F(\pi)-F(\pi_{t+1})
&\le \langle g_{t+1},\pi-\pi_{t+1}\rangle_\D\\
&\le\frac1{\eta_t}
\left\langle\nabla h(\pi_{t+1})-\nabla h(\pi_t),
\pi-\pi_{t+1}\right\rangle_\D\\
&=\frac1{\eta_t}
\left[
D(\pi\|\pi_t)-D(\pi\|\pi_{t+1})-D(\pi_{t+1}\|\pi_t)
\right].
\label{eq:prox-three-point}
\end{align}
For $t\ge1$, $\pi_t$ is feasible in the next subproblem. Comparing the optimum with $\pi_t$ gives
\begin{equation}
F(\pi_{t+1})-\frac{1}{\eta_t}D(\pi_{t+1}\|\pi_t)
\ge F(\pi_t),
\label{eq:prox-monotone}
\end{equation}
so the objective values are nondecreasing from $\pi_1$ onward.

Set $\pi=\pi^\star$ in \eqref{eq:prox-three-point}, multiply by $\eta_t$, and sum from $t=0$ to $T-1$:
\begin{equation}
\sum_{t=0}^{T-1}\eta_t
\big(F(\pi^\star)-F(\pi_{t+1})\big)
\le
D(\pi^\star\|\piref).
\label{eq:prox-telescope}
\end{equation}
Because the gaps are nonincreasing, the left-hand side is at least
$\big(\sum_t\eta_t\big)(F(\pi^\star)-F(\pi_T))$, which proves \eqref{eq:last-iterate-rate}.

\begin{corollary}[Convergence of the surplus profile]
\label{cor:surplus-convergence}
Under the conditions of Theorem~\ref{thm:last-iterate}, let
$s^\star=s(\pi^\star)$ denote the unique surplus vector induced by a
Nash-optimal policy. Then
\begin{equation}
\|s(\pi_T)-s^\star\|_2^2
\le
\frac{2D(\pi^\star\|\piref)}
{\sum_{t=0}^{T-1}\eta_t}.
\label{eq:last-iterate-surplus}
\end{equation}
Hence a constant proximal step gives
$\|s(\pi_T)-s^\star\|_2=O(T^{-1/2})$.
\end{corollary}

\begin{proof}
This is a stronger consequence of Nash-welfare convergence rather than a
separate requirement for the NBPO update. For every objective,
$V_k(\pi)\in[-1/2,1/2]$, and
$d_k=V_k(\piref)\le0$ because $\nu=\piref$ is feasible and
$J_k(\piref,\piref)=0$. Hence every positive surplus is at most one. On
$(0,1]^K$, the function
$G(s)=\sum_k\log s_k$ is $1$-strongly concave in Euclidean norm.
Proposition~\ref{prop:nash-outcome} implies that all Nash-optimal policies induce the
same game-value vector, and therefore the same surplus vector $s^\star$.
First-order optimality over the convex comprehensive surplus set induced by
\eqref{eq:comprehensive-hull} gives
\begin{equation}
G(s^\star)-G(s)
\ge
\frac12\|s-s^\star\|_2^2.
\label{eq:surplus-strong-concavity}
\end{equation}
Applying this inequality to $s=s(\pi_T)$ and using
\eqref{eq:last-iterate-rate} proves \eqref{eq:last-iterate-surplus}.
\end{proof}

\section{Objective-Wise Evaluation}
\label{app:objectivewise}
Table~\ref{tab:neutral-objectivewise} expands the worst and average scores in Table~\ref{tab:multi-dataset-external}. Every cell must be produced by the same prompt-disjoint evaluator protocol used for the corresponding main-table row; no test-selected policy or earlier checkpoint may be substituted.

% Optional generated replacement: objectivewise_generated.tex should contain
% the complete table* environment and preserve label tab:neutral-objectivewise.
\IfFileExists{\NBPOResultDir/objectivewise_generated.tex}{%
  \input{\NBPOResultDir/objectivewise_generated.tex}%
}{%
\begin{table*}[t]
\centering\scriptsize
\setlength{\tabcolsep}{2.4pt}
\caption{Objective-wise held-out win rates $W_k^{\rm eval}$ against the common reference. Entries are mean $\pm$ standard deviation over three training seeds. Higher is better.}
\label{tab:neutral-objectivewise}
\begin{tabular}{lcccccc}
\toprule
method & objective 1 & objective 2 & objective 3 & objective 4 & worst & average\\
\midrule
\multicolumn{7}{l}{\textit{\textsc{UltraFeedback}: instruction following, truthfulness, honesty, helpfulness}}\\
Base $=\piref$ & \pending & \pending & \pending & \pending & \pending & \pending \\
\NBPO{} (ours) & \pending & \pending & \pending & \pending & \pending & \pending \\
Fixed-reference Nash & \pending & \pending & \pending & \pending & \pending & \pending \\
BT-RM--Nash & \pending & \pending & \pending & \pending & \pending & \pending \\
Game-utilitarian & \pending & \pending & \pending & \pending & \pending & \pending \\
Game-KS & \pending & \pending & \pending & \pending & \pending & \pending \\
BT-RM--utilitarian & \pending & \pending & \pending & \pending & \pending & \pending \\
PROSPER/MaxEntBW & \pending & \pending & \pending & \pending & \pending & \pending \\
MOPO & \pending & \pending & \pending & \pending & \pending & \pending \\
RACO & \pending & \pending & \pending & \pending & \pending & \pending \\
\midrule
\multicolumn{7}{l}{\textit{\textsc{TL;DR}: coverage, faithfulness, conciseness, helpfulness}}\\
Base $=\piref$ & \pending & \pending & \pending & \pending & \pending & \pending \\
\NBPO{} (ours) & \pending & \pending & \pending & \pending & \pending & \pending \\
Fixed-reference Nash & \pending & \pending & \pending & \pending & \pending & \pending \\
BT-RM--Nash & \pending & \pending & \pending & \pending & \pending & \pending \\
Game-utilitarian & \pending & \pending & \pending & \pending & \pending & \pending \\
Game-KS & \pending & \pending & \pending & \pending & \pending & \pending \\
BT-RM--utilitarian & \pending & \pending & \pending & \pending & \pending & \pending \\
PROSPER/MaxEntBW & \pending & \pending & \pending & \pending & \pending & \pending \\
MOPO & \pending & \pending & \pending & \pending & \pending & \pending \\
RACO & \pending & \pending & \pending & \pending & \pending & \pending \\
\midrule
\multicolumn{7}{l}{\textit{\textsc{SafeRLHF}: helpfulness, harmlessness}}\\
Base $=\piref$ & \pending & \pending & -- & -- & \pending & \pending \\
\NBPO{} (ours) & \pending & \pending & -- & -- & \pending & \pending \\
Fixed-reference Nash & \pending & \pending & -- & -- & \pending & \pending \\
BT-RM--Nash & \pending & \pending & -- & -- & \pending & \pending \\
Game-utilitarian & \pending & \pending & -- & -- & \pending & \pending \\
Game-KS & \pending & \pending & -- & -- & \pending & \pending \\
BT-RM--utilitarian & \pending & \pending & -- & -- & \pending & \pending \\
PROSPER/MaxEntBW & \pending & \pending & -- & -- & \pending & \pending \\
MOPO & \pending & \pending & -- & -- & \pending & \pending \\
RACO & \pending & \pending & -- & -- & \pending & \pending \\
\bottomrule
\end{tabular}
\end{table*}%
}

\section{Implementation and Hyperparameters}
\label{app:hyperparameters}
The following tables distinguish settings fixed before experimentation from quantities selected on validation data or measured after execution. Exact model and dataset revisions, split hashes, rubric hashes, and run-registry identifiers must accompany the released artifacts.

\begin{table*}[t]
\centering\scriptsize
\setlength{\tabcolsep}{2.2pt}
\caption{Data generation, preference labeling, and evaluation protocol. A dash indicates a value to be inserted from the immutable run manifest.}
\label{tab:data-sampling-hparams}
\begin{tabular}{@{}p{0.26\textwidth}p{0.235\textwidth}p{0.235\textwidth}p{0.235\textwidth}@{}}
\toprule
setting & \textsc{UltraFeedback} & \textsc{TL;DR} & \textsc{SafeRLHF}\\
\midrule
Objective panel & instruction following; truthfulness; honesty; helpfulness & coverage; faithfulness; conciseness; helpfulness & helpfulness; harmlessness\\
Train / validation / test prompts & $7{,}000/500/1{,}000$ & $3{,}000/500/1{,}000$ & $3{,}000/500/1{,}000$\\
Dataset revision and split hash & \pending & \pending & \pending\\
Overlap/decontamination audit & exact + approximate benchmark filtering; \pending removed & source-group split; \pending removed & source-group split; \pending removed\\
Base/reference policy & \multicolumn{3}{c}{\texttt{Llama-3.1-8B-Instruct}; exact revision \pending}\\
Training preference judge & \multicolumn{3}{c}{\texttt{Qwen3-32B}, thinking disabled, temperature $0$, two presentation orders}\\
Primary / secondary evaluators & \multicolumn{3}{c}{\texttt{Llama-3.3-70B-Instruct} / \texttt{Phi-4}; exact revisions \pending}\\
Learner / comparator responses per prompt & $4/4$ & $4/4$ & $4/4$\\
Cross pairs per prompt and objective & $16$ & $16$ & $16$\\
Training response generation & temp. $0.9$, top-$p$ $0.95$, $512$ tokens & temp. $0.9$, top-$p$ $0.95$, $384$ tokens & temp. $0.9$, top-$p$ $0.95$, $512$ tokens\\
Evaluation response generation & \multicolumn{3}{c}{temp. $0.7$, top-$p$ $0.9$, at most $1{,}024$ new tokens}\\
Cycle-audit prompts & $1{,}000$ & $500$ & $500$\\
Unique semantic comparisons & \pending & \pending & \pending\\
Order-swapped judge calls incl. retries & \pending & \pending & \pending\\
Invalid / retry / tie rates & \pending & \pending & \pending\\
Rubric, prompt, and judge hashes & \pending & \pending & \pending\\
\bottomrule
\end{tabular}
\end{table*}

\begin{table*}[t]
\centering\scriptsize
\setlength{\tabcolsep}{2.2pt}
\caption{Common neural policy-training protocol. Validation-selected values are left fill-ready and must be identical across matched methods unless the row explicitly states otherwise.}
\label{tab:policy-training-hparams}
\begin{tabular}{@{}p{0.26\textwidth}p{0.235\textwidth}p{0.235\textwidth}p{0.235\textwidth}@{}}
\toprule
setting & \textsc{UltraFeedback} & \textsc{TL;DR} & \textsc{SafeRLHF}\\
\midrule
Adaptation / precision & full fine-tuning / bf16 & full fine-tuning / bf16 & full fine-tuning / bf16\\
Optimizer / scheduler & Adafactor / cosine & Adafactor / cosine & Adafactor / cosine\\
Warmup / weight decay / grad. clip & $0.10/0/1.0$ & $0.10/0/1.0$ & $0.10/0/1.0$\\
Maximum sequence length & $2{,}048$ & $4{,}096$ & $2{,}048$\\
Microbatch / accumulation / global batch & $2/16/32$ & $1/32/32$ & $2/16/32$\\
Learning-rate candidates; selected & $\{2,5,10\}\!\times\!10^{-7}$; \pending & $\{2,5,10\}\!\times\!10^{-7}$; \pending & $\{2,5,10\}\!\times\!10^{-7}$; \pending\\
Policy updates candidates; selected & $\{300,600\}$; \pending & $\{300,600\}$; \pending & $\{300,600\}$; \pending\\
Proximal $\eta$ candidates; selected & $\{0.125,0.25,0.5,1\}$; \pending & $\{0.125,0.25,0.5,1\}$; \pending & $\{0.125,0.25,0.5,1\}$; \pending\\
Opponent $\beta$ candidates; selected & $\{0.15,0.25,0.5,1\}$; \pending & $\{0.15,0.25,0.5,1\}$; \pending & $\{0.15,0.25,0.5,1\}$; \pending\\
Fixed-point $R$ candidates; selected & $\{1,3,5\}$; \pending & $\{1,3,5\}$; \pending & $\{1,3,5\}$; \pending\\
Dual solver / tolerance / box & \pending & \pending & \pending\\
Regression target & \multicolumn{3}{c}{Rao--Blackwellized soft pairwise target; sampled Bernoulli target used only as an ablation}\\
Validation selection & \multicolumn{3}{c}{lexicographic: maximize worst objective, then average within $0.005$, then choose smaller reference KL}\\
Tuning / final seeds & \multicolumn{3}{c}{$11$ / $\{42,43,44\}$}\\
Hardware & \multicolumn{3}{c}{$1\times$ H200 per 8B training run; $4\times$ H200 for tensor-parallel judging}\\
Wall-clock / GPU-hours per run & \pending & \pending & \pending\\
Source commit / frozen-config hash & \pending & \pending & \pending\\
\bottomrule
\end{tabular}
\end{table*}

\section{Reward-Model Baseline Protocol}
\label{app:rm-baselines}
BT-RM--Nash and BT-RM--utilitarian share the same objective-specific reward models, soft Bradley--Terry labels, reference normalization, response pool, and policy-fitting budget. They differ only in cross-objective aggregation.

\begin{table*}[t]
\centering\scriptsize
\setlength{\tabcolsep}{2.2pt}
\caption{Training and validation record for the matched objective-specific Bradley--Terry reward models.}
\label{tab:rm-baseline-protocol}
\begin{tabular}{@{}p{0.26\textwidth}p{0.235\textwidth}p{0.235\textwidth}p{0.235\textwidth}@{}}
\toprule
setting & \textsc{UltraFeedback} & \textsc{TL;DR} & \textsc{SafeRLHF}\\
\midrule
Number of objective-specific RMs & $4$ & $4$ & $2$\\
Backbone / scalar head & \multicolumn{3}{c}{\texttt{Llama-3.1-8B-Instruct}; scalar head on final non-padding response token}\\
Training labels / loss & \multicolumn{3}{c}{swap-averaged soft preference probabilities / soft Bradley--Terry loss}\\
Training semantic pairs & \pending & \pending & \pending\\
Backbone LR candidates; selected & $\{5\!\times\!10^{-7},10^{-6},2\!\times\!10^{-6}\}$; \pending & same; \pending & same; \pending\\
Head LR / epochs / global batch & \pending & \pending & \pending\\
Selected validation soft NLL by objective & \pending & \pending & \pending\\
Hard accuracy / Brier / ECE by objective & \pending & \pending & \pending\\
Reference reward mean/std and hash & \pending & \pending & \pending\\
BT-RM--Nash policy config & \pending & \pending & \pending\\
BT-RM--utilitarian policy config & \pending & \pending & \pending\\
RM training + policy GPU-hours & \pending & \pending & \pending\\
\bottomrule
\end{tabular}
\end{table*}

\section{External-Baseline Selection and Compute}
\label{app:external-baselines}
External methods retain their defining optimization rule but use the shared data and evaluation protocol. Table~\ref{tab:external-baseline-protocol} records the exact implementation revision, validation search, selected operating point, and compute so that a family-valued baseline cannot benefit from test-set selection.

\begin{table*}[t]
\centering\scriptsize
\setlength{\tabcolsep}{2.8pt}
\caption{Implementation and validation record for external multi-objective baselines.}
\label{tab:external-baseline-protocol}
\begin{tabular}{p{2.4cm}p{3.0cm}p{3.3cm}p{3.1cm}p{2.0cm}}
\toprule
method & implementation / revision & validation search & selected operating point & GPU-hours\\
\midrule
PROSPER/MaxEntBW & \pending & method-specific parameters: \pending & \pending & \pending\\
MOPO & \pending & each objective as primary; relative threshold $\alpha\in\{0.25,0.5,0.75\}$ & \pending & \pending\\
RACO & \pending & $\beta_{\rm DPO}\in\{0.1,0.2\}$, $c\in\{0.2,0.4,0.6\}$, predeclared weight vectors & \pending & \pending\\
\bottomrule
\end{tabular}
\end{table*}

\section{Additional External Evaluation}
\label{app:external-evaluation}
Safety transfer uses the same frozen \textsc{UltraFeedback}-trained checkpoints as Table~\ref{tab:general-benchmark-completion}. Because attack-success rate has the opposite direction from the main metrics, it is reported separately.

\begin{table}[t]
\centering\footnotesize
\setlength{\tabcolsep}{5.0pt}
\caption{HarmBench attack-success rate (ASR; lower is better). Entries are mean $\pm$ standard deviation over the same three training seeds.}
\label{tab:safety-transfer}
\begin{tabular}{lc}
\toprule
method & HarmBench ASR $\downarrow$\\
\midrule
Base $=\piref$ & \pending\\
\NBPO{} (ours) & \pending\\
Fixed-reference Nash & \pending\\
BT-RM--Nash & \pending\\
Game-utilitarian & \pending\\
Game-KS & \pending\\
BT-RM--utilitarian & \pending\\
PROSPER/MaxEntBW & \pending\\
MOPO & \pending\\
RACO & \pending\\
\bottomrule
\end{tabular}
\end{table}

\section{Optimization Diagnostics and Robustness}
\label{app:optimization-diagnostics}
The following diagnostics separate finite-pool optimization from neural realization and test the conditions under which adaptive opponents should matter. All subset definitions and hyperparameter candidates are fixed before final policy evaluation.

\begin{table*}[t]
\centering\scriptsize
\setlength{\tabcolsep}{3.0pt}
\caption{Validation-stage selection record. The selected setting is frozen before the three final training seeds.}
\label{tab:optimization-ablation}
\begin{tabular}{lp{4.0cm}ccccc}
\toprule
factor & candidates & selected & worst $W^{\rm eval}$ & avg. $W^{\rm eval}$ & norm. reg. MSE & fixed-point residual\\
\midrule
Opponent temperature $\beta$ & $\{0.15,0.25,0.5,1\}$ & \pending & \pending & \pending & \pending & \pending\\
Fixed-point budget $R$ & $\{1,3,5\}$ & \pending & \pending & \pending & \pending & \pending\\
Proximal coefficient $\eta$ & $\{0.125,0.25,0.5,1\}$ & \pending & \pending & \pending & \pending & \pending\\
Policy learning rate & $\{2,5,10\}\times10^{-7}$ & \pending & \pending & \pending & \pending & \pending\\
Policy updates & $\{300,600\}$ & \pending & \pending & \pending & \pending & \pending\\
\bottomrule
\end{tabular}
\end{table*}

\begin{table*}[t]
\centering\scriptsize
\setlength{\tabcolsep}{2.4pt}
\caption{Population-target and neural-realization diagnostics on the frozen \textsc{UltraFeedback} pool. Exact scores are computed from the finite-pool solution; trained scores use the final neural policy.}
\label{tab:finite-pool-realization}
\resizebox{\textwidth}{!}{%
\begin{tabular}{lccccccccc}
\toprule
method & exact min $s$ & exact avg. $s$ & trained min $s$ & trained avg. $s$ & min-$s$ gap & dual/KKT residual & norm. reg. MSE & target corr. & sign agreement\\
\midrule
\NBPO{} & \pending & \pending & \pending & \pending & \pending & \pending & \pending & \pending & \pending\\
Fixed-reference Nash & \pending & \pending & \pending & \pending & \pending & \pending & \pending & \pending & \pending\\
BT-RM--Nash & \pending & \pending & \pending & \pending & \pending & \pending & \pending & \pending & \pending\\
Game-utilitarian & \pending & \pending & \pending & \pending & \pending & \pending & \pending & \pending & \pending\\
Game-KS & \pending & \pending & \pending & \pending & \pending & \pending & \pending & \pending & \pending\\
BT-RM--utilitarian & \pending & \pending & \pending & \pending & \pending & \pending & \pending & \pending & \pending\\
\bottomrule
\end{tabular}}
\end{table*}

\begin{table*}[t]
\centering\scriptsize
\setlength{\tabcolsep}{3.0pt}
\caption{Pairwise update-score sign flips (\%) under three independently rescaled objective coordinate systems. Lower is better; zero denotes exact numerical invariance.}
\label{tab:uf-scale-covariance}
\begin{tabular}{lcccc}
\toprule
method & $c^{(A)}$ & $c^{(B)}$ & $c^{(C)}$ & mean\\
\midrule
\NBPO{} & \pending & \pending & \pending & \pending\\
Fixed-reference Nash & \pending & \pending & \pending & \pending\\
BT-RM--Nash & \pending & \pending & \pending & \pending\\
Game-utilitarian & \pending & \pending & \pending & \pending\\
Game-KS & \pending & \pending & \pending & \pending\\
BT-RM--utilitarian & \pending & \pending & \pending & \pending\\
PROSPER/MaxEntBW & \pending & \pending & \pending & \pending\\
MOPO & \pending & \pending & \pending & \pending\\
RACO & \pending & \pending & \pending & \pending\\
\bottomrule
\end{tabular}
\end{table*}

\begin{table*}[t]
\centering\scriptsize
\setlength{\tabcolsep}{2.5pt}
\caption{Cycle-conditioned representation comparison. Prompt subsets are defined from the frozen preference matrices before policy evaluation. Higher $W^{\rm eval}$ is better.}
\label{tab:cycle-conditioned}
\resizebox{\textwidth}{!}{%
\begin{tabular}{lccccccccc}
\toprule
subset & prompts & hard-cycle rate & soft-cycle mass & BT deviance & \NBPO{} worst $W$ & fixed-ref. worst $W$ & BT-RM worst $W$ & \NBPO{}$-$fixed & \NBPO{}$-$BT-RM\\
\midrule
All audited prompts & \pending & \pending & \pending & \pending & \pending & \pending & \pending & \pending & \pending\\
Low-intransitivity quartile & \pending & \pending & \pending & \pending & \pending & \pending & \pending & \pending & \pending\\
High-intransitivity quartile & \pending & \pending & \pending & \pending & \pending & \pending & \pending & \pending & \pending\\
Low-BT-deviance quartile & \pending & \pending & \pending & \pending & \pending & \pending & \pending & \pending & \pending\\
High-BT-deviance quartile & \pending & \pending & \pending & \pending & \pending & \pending & \pending & \pending & \pending\\
\bottomrule
\end{tabular}}
\end{table*}

\section{Judge Instantiation and Preference Labeling}
\label{app:labeling}
\paragraph{Training judge.}
For each objective, \texttt{Qwen3-32B} compares two responses under an objective-specific rubric with reasoning traces disabled. Each semantic pair is queried in both presentation orders. Verdicts are restricted to \texttt{[[A]]}, \texttt{[[B]]}, or \texttt{[[TIE]]}; scoring them as $1$, $0$, or $1/2$ and averaging the two orders produces an exactly skew-symmetric empirical margin matrix.

\paragraph{Objective panels.}
\textsc{UltraFeedback} uses instruction following, truthfulness, honesty (calibration and non-deception), and helpfulness. \textsc{TL;DR} uses coverage, faithfulness, conciseness, and helpfulness. \textsc{SafeRLHF} uses helpfulness and harmlessness. The judge is instructed to evaluate only the named objective; exact rubric text and hashes are included in the released configuration.

\paragraph{Pairwise template.}
The common interface is
\begin{quote}\small\ttfamily\raggedright\sloppy
System: You are an impartial evaluator comparing two assistant responses. Evaluate ONLY the objective stated below. Choose A only when A is materially better, B only when B is materially better, and TIE otherwise. Return exactly [[A]], [[B]], or [[TIE]].

User: [OBJECTIVE] \{objective name and rubric\} [USER REQUEST / SOURCE] \{prompt\} [RESPONSE A] \{A\} [RESPONSE B] \{B\} [VERDICT]
\end{quote}
The second call swaps responses A and B while leaving every other field unchanged.

\paragraph{From verdicts to game margins.}
If $q_{AB}\in\{1,1/2,0\}$ scores the first call from the perspective of original response $A$, and $q_{BA}$ scores the swapped call from the perspective of original response $B$, then
\[
\widehat P_k(A\succ B\mid x)
=\tfrac12\bigl(q_{AB}+1-q_{BA}\bigr),
\qquad
\widehat\Delta_k(A,B\mid x)
=\widehat P_k(A\succ B\mid x)-\tfrac12.
\]
The shared probability is used directly in the soft Bradley--Terry loss. For \NBPO{}, reference responses are reweighted by the regularized-opponent expression in \eqref{eq:regularized-opponent}, and pairwise differences of the resulting response-level scores form the Rao--Blackwellized regression target corresponding to \eqref{eq:binary-target-mean}.


\end{document}